\documentclass[12pt,a4paper]{article}

\usepackage[T1]{fontenc}
\usepackage[utf8]{inputenc}
\usepackage{lmodern}
\usepackage{amsmath,amssymb,bm}
\usepackage{graphicx}
\usepackage{array}
\usepackage{booktabs}
\usepackage[a4paper,margin=2.5cm]{geometry}
\usepackage[hidelinks]{hyperref}
\usepackage{setspace}

\setstretch{1.15}
\setlength{\parindent}{0.7cm}
\setlength{\parskip}{0pt}
\setlength{\emergencystretch}{2em}
\newcolumntype{L}[1]{>{\raggedright\arraybackslash}p{#1}}

\title{Refractive Gravity: a phase--solid effective field theory\\
with a regular horizonless compact branch}
\author{George Lotishvili\\
\small East European University, Tbilisi, Georgia\\
\small \texttt{g\_lotishvili@eeu.edu.ge}}
\date{}

\begin{document}

\maketitle

\begin{abstract}
We present Refractive Gravity (RefG), a metric effective field theory
consisting of Einstein--Hilbert gravity, one shift-symmetric phase scalar, and
three solid scalars, in the supersolid class of self-gravitating media. The
scalar sector is a minimal low-order polynomial response basis in the standard
medium invariants; the tensor sector lies in the constant-$G_4$ Horndeski
class, so gravitational waves propagate at the speed of light identically. In
the weak static regime, the medium stress reproduces Newton's law through a
refractive source--index chain. The Solar branch coincides with general
relativity through the second post-Newtonian order, yielding $\gamma=\beta=1$,
a vanishing preferred-frame source combination at the displayed order, and a
light-propagation coefficient $q_{2{\rm PN}}=7/4$. In the strong-field regime,
we show that the projected phase-pressure deficit of the same scalar sector
sources the exponential (Papapetrou) exterior exactly. It reproduces the
Einstein tensor component by component without a propagating ghost --- a geometry
previously realized in general relativity only by a phantom scalar. A
$C^2$-matched regular core removes the wormhole throat, rendering the
geometry horizonless and geodesically complete. The deficit-feedback
exponent is fixed by the volume readout of the asymptotic mass, leaving a
parameter-free realizing window $r_s/2<r_c<r_s$. The resulting static branch
is falsifiable: the shadow radius exceeds the Schwarzschild value by
$4.6\%$, within current Sagittarius~A* bounds at the $2\sigma$ level, and the
ISCO frequency is $0.931$ of the general-relativity value at the same mass.
\end{abstract}

\section{Introduction}
\label{sec:intro}

General relativity (GR) has passed every weak-field and radiative test to
date, from Solar-System post-Newtonian measurements to the direct detection of
gravitational waves \cite{Will2014,Bertotti2003,Abbott2017GW}. The remaining
frontier is the strong, static gravitational field. Horizon-scale images of M87* and
Sagittarius~A* \cite{EHT2019M87,EHT2022SgrA} and the ringdown phase of binary
mergers now probe the geometry within a few Schwarzschild radii. In this
regime the black-hole paradigm faces its sharpest alternatives: horizonless
ultracompact objects whose exteriors mimic a black hole while their interiors
remain regular \cite{CardosoPani2019}. Quantitative tests of this kind
require concrete theories: actions whose field equations select a definite
exterior, complementing parametrized metric approaches
\cite{CliftonEtAl2012}.

A systematic source of such theories is the effective field theory (EFT) of
gravitating media. Depending on the symmetries imposed on a set of scalar
fields, one obtains superfluids, solids, framids, or supersolids
\cite{NicolisEtAl2015,CeloriaComelliPilo2017}. The same construction underlies
solid inflation \cite{Endlich2013} and the EFT of dark energy
\cite{Cheung2008,Gubitosi2013,BelliniSawicki2014}. Refractive Gravity (RefG)
is a theory within this class, comprising one shift-symmetric phase scalar
$\Phi$ and three solid scalars $\phi^A$ coupled minimally to the metric --- the
field content
of a supersolid medium. The word ``medium'' is used throughout as the physical
reading of this scalar sector, not as a microscopic substrate assumed in
advance; matter couples only to the metric, and every statement below follows
from the action, metric variation, and branch-specific field equations. The
name of the theory refers to its weak-field mechanism: the static balance of
the medium stress defines an effective refractive profile for slow matter and
light. Newton's law then emerges as the weak spherical limit of this
source--index chain.

The strong-field analysis in this paper concerns a geometry with a long
history: the exponential metric
$ds^2=e^{-r_s/r}c^2dt^2-e^{r_s/r}(dr^2+r^2d\Omega^2)$. This metric was first
written by Papapetrou \cite{Papapetrou1954}, later served as the basis of an
alternative field-equation program whose status remained contested
\cite{Yilmaz1958,Misner1999,AlleyYilmaz1995}, and was recently re-examined within
GR as a traversable wormhole by Boonserm, Ngampitipan, Simpson, and Visser
\cite{Boonserm2018}. The test-particle and wave phenomenology of this exterior
is known: the photon sphere, critical impact parameter, ISCO radius, and
Regge--Wheeler potentials were derived in \cite{Boonserm2018}, the shadow size
and ISCO frequency shift in \cite{MakukovMychelkin2018}, and the confrontation
with Event Horizon Telescope data in \cite{Vagnozzi2023}. None of the
exterior geodesic numbers quoted in this paper are new; they are used with
attribution and a one-line normalization table.

What has remained unsettled is the physical source of this geometry. Within
GR, the exponential exterior requires an effective stress that violates the
radial null energy condition at every radius, and its canonical realization is
a massless ghost (phantom) scalar \cite{Boonserm2018}. Alternative realizations
in the literature --- antiscalar backgrounds \cite{MakukovMychelkin2018} or
modified field-equation programs \cite{Yilmaz1958} --- either involve a ghost
degree of freedom or operate outside the standard field-equation framework. A
ghost-free field-theoretic source for this geometry, with a regular interior in
place of the wormhole throat, has been missing.

In this paper we supply such a source, and we delimit exactly what it does
and does not establish. We show three results. First, the projected
phase-pressure deficit of the RefG scalar sector --- a single higher-sector
operator with no propagating time-kinetic degree of freedom --- sources the
exponential exterior exactly. Its stress reproduces the Einstein tensor in all
four diagonal components with one common profile function --- a four-component
functional match that no constant rescaling could imitate. Second, the
finite-core load selector, which decides between the weak (Solar) branch and
the compact branch, closes into a parameter-free statement: the feedback
exponent is fixed by the volume readout of the asymptotic mass, narrowing the
realizing window to $r_s/2<r_c<r_s$ in the core radius --- equivalently, a
compactness band with no adjustable constant. Third, a $C^2$-matched
logarithmic core replaces the region containing the wormhole throat. The
resulting geometry is horizonless, regular at the center, and geodesically
complete, with an effective source that is bounded and continuous at the
junction.

On this static branch, the observational markers are inherited from the known
exterior and remain fixed: the shadow radius exceeds the Schwarzschild value by
$4.6\%$ and the ISCO frequency is $0.931$ of the GR value at the same mass.
In contrast to model families with a free closeness parameter, which can
approach the Kerr limit arbitrarily closely \cite{CardosoPani2019}, these
numbers cannot be adjusted: current Sagittarius~A*
measurements already place the shadow prediction under mild tension at the
$1\sigma$ level while allowing it at $2\sigma$, and next-generation
horizon-scale instruments can falsify the branch. We also confront
the generic structural risk of this object class: the cored geometry
possesses the stable inner light ring required by the light-ring pairing
theorem \cite{CunhaBertiHerdeiro2017}. We compute its location and the depth
of the trapping well, which proves to be shallow across the whole realizing
window; the nonlinear fate of trapped modes remains open, as it does for the
entire horizonless class \cite{CunhaEtAl2023}.

The scope of the paper is fixed as follows. All strong-field statements
refer to the static, spherically symmetric branch. The rotating exterior, the
quasinormal-mode and echo spectrum, a matter equation of state for the core,
and the collapse dynamics that would select the branch for a given
astrophysical object are open problems; they are collected in the Discussion,
alongside the specific reason each one does not affect the static results. The
weak-field sector of RefG coincides with GR through the second post-Newtonian
order; those results are consistency checks, and the only GR-discriminating
predictions of the paper are the compact-branch markers.

The paper is organized as follows. Section~\ref{sec:theory} introduces the
field content, the design requirements, the action, and the branch structure,
and states the tensor-sector and cosmological-background properties.
Section~\ref{sec:weak} derives the refractive source--index chain, Newton's
law, and the Solar post-Newtonian results. Section~\ref{sec:stability}
presents the local stability analysis of the scalar-longitudinal sector.
Section~\ref{sec:compact} contains the compact branch: the exterior, the
raw-invariant obstruction, the phase-normalized action and the deficit-source
realizability theorem, the parameter-free selector, the regular core, and the
geodesic markers. Section~\ref{sec:discussion} confronts the branch with
observations and states the open problems and falsifiability conditions.
Section~\ref{sec:conclusion} concludes. Two appendices collect the stability
diagnostics and the Solar-system derivations.

\section{Fields, action, and branch structure}
\label{sec:theory}

\subsection{Design requirements}
\label{sec:design}

The theory is constructed to satisfy five requirements, each mapping to
a specific structural choice below:
\begin{enumerate}
\item exact recovery of Newton's law and of GR light bending in the weak
static limit (the refractive source--index chain of
Section~\ref{sec:weak});
\item $\gamma=\beta=1$ and a vanishing preferred-frame source combination on
the Solar branch (the coefficient family of Section~\ref{sec:solar});
\item luminal tensor waves, $c_g=c$, identically rather than by tuning (the
constant-$G_4$ Horndeski projection of Section~\ref{sec:tensor});
\item a ghost-free and gradient-stable scalar sector around the relevant
backgrounds (the completed kinetic window of Section~\ref{sec:stability});
\item the existence of a static compact branch with a regular, horizonless
exterior realized by the same scalar sector (Section~\ref{sec:compact}).
\end{enumerate}
The action is therefore presented as the conclusion of these requirements,
not as a postulate.

\subsection{Field content and invariants}
\label{sec:fields}

The gravitational variables are the metric $g_{\mu\nu}$ (signature $(+---)$),
one phase scalar $\Phi$ with internal shift symmetry
$\Phi\to\Phi+{\rm const}$, and three solid scalars $\phi^A$, $A=1,2,3$, with
internal shifts and rotations $\phi^A\to R^A{}_B\phi^B+a^A$. This is the field
content of a supersolid medium in the classification of gravitating media
\cite{NicolisEtAl2015,CeloriaComelliPilo2017}. The lowest-derivative
invariants consistent with these symmetries are
\begin{align*}
Y &= g^{\mu\nu}\partial_\mu\Phi\,\partial_\nu\Phi,\\
B^{AB} &= -g^{\mu\nu}\partial_\mu\phi^A\,\partial_\nu\phi^B,\\
I_1 &= {\rm Tr}\,B,\qquad
I_2=\frac{1}{2}\left[({\rm Tr}\,B)^2-{\rm Tr}(B^2)\right],\qquad
I_3=\det B.
\end{align*}
On the homogeneous background $\Phi=t$, $\phi^A=x^A$, one has $Y=1$,
$B^{AB}=\delta^{AB}$. In the convention $\mathcal L=M_*^4F$, the invariants
are dimensionless, the coefficients introduced below are medium response
parameters, and the density scale is carried by $M_*^4$. For a static,
spherical solid configuration, $B^{AB}$ has one radial and two tangential
eigenvalues, denoted $\lambda_r$ and $\lambda_t$; these eigenvalues appear
throughout the analysis of the compact branch.

\subsection{Action}
\label{sec:action}

The minimal low-order polynomial response basis in these invariants --- seven
terms, with a single phase--solid mixing term --- is given by
\begin{align*}
F_{\rm min}={}&
c_Y Y+c_{Y2}Y^2+c_{I1}I_1+c_{I1sq}I_1^2+c_{I2}I_2+c_{I3}I_3
+c_{YI1}YI_1.
\end{align*}
The full action is
\begin{align*}
S=\int d^4x\sqrt{-g}\left[
\frac{M_{\rm Pl}^2}{2}R+M_*^4 F_{\rm min}^{(H)}
+\mathcal L_{\rm EFT}^{>{\rm min}}
\right]+S_{\rm m}[g_{\mu\nu},\psi],
\end{align*}
where matter couples minimally to the metric alone. The term
$\mathcal L_{\rm EFT}^{>{\rm min}}$ denotes the higher EFT sector of the same
medium (one operator of which becomes active on the compact branch, as
detailed in Section~\ref{sec:compact}), and $F_{\rm min}^{(H)}$ is the same
polynomial evaluated on the phase-normalized strain invariants defined below.
The active stress of the scalar sector is obtained by metric variation,
\[
\Theta^{\rm RefG}_{\mu\nu}
=2\frac{\partial\mathcal L_{\rm RefG}}{\partial g^{\mu\nu}}
-g_{\mu\nu}\mathcal L_{\rm RefG},
\qquad
G_{\mu\nu}=8\pi G\,\Theta^{\rm RefG}_{\mu\nu}
\quad (T^{\rm m}_{\mu\nu}=0),
\]
with the overall sign appropriate to the $(+---)$ signature (for a canonical
scalar, this returns
$\Theta_{\mu\nu}=\partial_\mu\phi\,\partial_\nu\phi-g_{\mu\nu}\mathcal L$).
This stress is the source of every exterior considered below.

\subsection{Branch frames and phase normalization}
\label{sec:branches}

Let $H$ be a phase-normalization variable, and define the rescaled strain
invariants
\[
\widehat Y=e^{-2H}Y,\qquad
\widehat\lambda_r=e^{2H}\lambda_r,\qquad
\widehat\lambda_t=e^{2H}\lambda_t,
\]
where the hatted invariants $\widehat I_n$ are formed from
$\widehat\lambda_r,\widehat\lambda_t$ in the same manner that $I_n$ are formed
from $\lambda_r,\lambda_t$. Because no derivative of $H$ appears in
$F_{\rm min}^{(H)}$, the variable $H$ is auxiliary, and its variation gives the
algebraic constraint
\[
S_H\equiv\frac{\partial F_{\rm min}^{(H)}}{\partial H}=0 .
\]
A branch is a background regime of this single action. The weak-Solar branch
is the unloaded normalization $H=0$, on which $F_{\rm min}^{(H)}$ reduces to
the raw polynomial and the constraint $S_H=0$ is satisfied through second
post-Newtonian order by the Solar strain itself (the expansion is collected
in Appendix~\ref{app:solar}). The compact branch is the loaded normalization
$H=h$, where $h$ is the metric phase of the compact exterior. In this regime
the hatted invariants equal one, and because the Solar coefficient family is
tadpole-free at the unit state,
\[
F_{\rm min}(1,1,1)=0,\qquad
\left.\frac{\partial F_{\rm min}}{\partial\widehat Y}\right|_1=
\left.\frac{\partial F_{\rm min}}{\partial\widehat\lambda_r}\right|_1=
\left.\frac{\partial F_{\rm min}}{\partial\widehat\lambda_t}\right|_1=0,
\]
the polynomial sector carries zero stress and zero $\phi^A$ Euler residual on
the compact exterior, and $S_H=0$ holds exactly. The motivation for this
normalization is structural rather than aesthetic:
Section~\ref{sec:obstruction} exhibits an explicit tensor obstruction that
forbids the insertion of raw invariants on the compact branch, and the
phase-normalized frame resolves it. A global identification of $H$ with the solid determinant,
$H=-\ln(\lambda_r\lambda_t^2)/6$, is rejected by the same algebra: it
generates a nonzero weak-Solar stress at first order and admits no exact-GR
second-order solution (Appendix~\ref{app:solar}).

Four distinct symbols are used and never interchanged: $h_{\rm eff}$ is the
weak static source profile, $h$ is the compact metric phase, $H$ normalizes
the strain invariants, and $H_\Delta$ (introduced in
Section~\ref{sec:compact}) is the independent deficit-source scalar.

\subsection{Tensor sector}
\label{sec:tensor}

Under the pure phase projection, with $Y=-2X$, the theory lies in the minimal
Horndeski class \cite{Horndeski1974}, yielding
\[
G_2=-2c_YX+4c_{Y2}X^2,\qquad G_3=0,\qquad
G_4=\frac{M_{\rm Pl}^2}{2},\qquad G_5=0 .
\]
Because $G_4$ is constant and $G_5=0$, the tensor speed excess vanishes
identically, $\alpha_T=0$ and $c_g=c$
\cite{Creminelli2017}, consistent with the binary-neutron-star bound
$|c_g/c-1|\lesssim{\rm few}\times10^{-15}$
\cite{Abbott2017GW,Abbott2017GRB}. The solid sector is static-silent for
transverse-traceless (TT) perturbations on the homogeneous background used
here: its quadratic TT corrections to both the kinetic and gradient terms
vanish; consequently, the local dispersion relation is $\omega^2=c^2k^2$, and
the tensor mode remains massless on the Minkowski branch. Although solid and
supersolid media can generate a tensor mass on cosmological backgrounds
\cite{CeloriaComelliPilo2017}, the corresponding FLRW mass term for RefG
remains an open problem of the cosmological extension, and it does not affect
the local statements of this paper.

\subsection{Cosmological background}
\label{sec:cosmo}

The same scalar sector admits a spatially flat FLRW background on the
zero-current branch of the phase equation, on which the medium contributes
$\rho=-p=c_Y^2/(4c_{Y2})$ --- an exact $w=-1$ component at the background
level. This is an algebraic statement about the background equations; the
perturbation analysis, the early-time scaling bounds, and the confrontation
with cosmological data are deferred to a separate study. We record it here
because
a new gravitational EFT must at minimum admit a sensible cosmological
background, and this one does.

\subsection{Relation to known theories}
\label{sec:relation}

RefG sits at a definite point of the known theory space. Its field content
and symmetries are those of a supersolid \cite{NicolisEtAl2015}; its
self-gravitating dynamics parallels those of the media of
\cite{CeloriaComelliPilo2017}; and its phase projection is minimal Horndeski
with a constant $G_4$. The timelike gradient of $\Phi$ plays a role analogous
to that of the aether vector in Einstein-aether theory, with the distinction
that here it is the gradient of a shift-symmetric scalar, while the solid
sector supplies the spatial frame. The compact exterior connects the theory
to the exponential-metric literature cited in the Introduction, but replaces
the source mechanism: where GR realizations require a phantom scalar
\cite{Boonserm2018}, RefG uses a projected deficit operator with no
propagating time-kinetic degree of freedom (Section~\ref{sec:compact}).

\section{Weak field: the refractive chain and the Solar system}
\label{sec:weak}

\subsection{Source--index chain and Newton's law}
\label{sec:newton}

The weak static metric is written as $g_{tt}=\exp(-2h_{\rm eff})$. For slow
test matter, the corresponding source index is $n_{\rm eff}=\exp(h_{\rm eff})$,
the effective potential per unit mass is $V_{\rm eff}/m=-c^2h_{\rm eff}$, and
the radial acceleration is given by
\[
a_r=c^2\frac{d\ln n_{\rm eff}}{dr}=c^2h_{\rm eff}' .
\]
The profile $h_{\rm eff}$ is sourced by the active medium stress. For a
static, spherical configuration with anisotropic pressure contrast
$\Delta_p=p_{\rm tan}-p_{\rm rad}$, the Bianchi/Tolman--Oppenheimer--Volkoff
balance of that stress reads
\[
p_{\rm rad}'+\rho_{\rm eff}\Phi_N'-\frac{2\Delta_p}{r}=0,
\qquad \Phi_N=-c^2h_{\rm eff},
\]
which locates the active pressure gradient in the refractive profile. The
deficit component of the same source is quadratic in the refractive gradient,
\[
\Pi_\Delta=\frac{e^{-2h_{\rm eff}}(h_{\rm eff}')^2}{8\pi G},
\]
and is the weak-field shadow of the operator that will carry the compact
branch. In the local spherical weak limit this term is subleading, $\Delta_p$
is absent at leading order, and the exterior source-free equation is
\[
\left(r^2h_{\rm eff}'\right)'=0 .
\]
Imposing Gauss normalization on the localized source charge fixes
$-r^2h_{\rm eff}'=GM/c^2$, yielding
\[
h_{\rm eff}(r)=\frac{GM}{c^2r},\qquad
\Phi_N=-\frac{GM}{r},\qquad
a_r=-\frac{GM}{r^2}.
\]
Consequently, Newton's law emerges as the weak spherical limit of the chain
$\Theta^{\rm RefG}_{\mu\nu}\to h_{\rm eff}\to n_{\rm eff}\to a_r$.

\subsection{Solar post-Newtonian branch}
\label{sec:solar}

The Solar calculation is performed in a Schwarzschild-like areal coordinate
with $u=GM/(c^2R)$,
\begin{align*}
g_{tt}^{A}&=1-2u+2(\beta_A-1)u^2+\cdots,\\
g_{RR}^{A}&=-(1+2\gamma u+a_2u^2+\cdots),
\end{align*}
which relates to the standard isotropic PPN metric via
$\beta_{\rm PPN}=\beta_A+\gamma-1$. Solving the first-order field equations
yields $\gamma=1$; at second order the medium-strain solution fixes $a_2=4$
(Appendix~\ref{app:solar}), and the same branch then gives $\beta_A=1$, hence
\[
\gamma=1,\qquad \beta_{\rm PPN}=1,\qquad b_2=\frac{a_2-\gamma^2}{2}=\frac32 .
\]
In the static scalar reduction, the preferred-frame source combination is
\[
K_{\rm pf}
=-c_{I1}-6c_{I1sq}-2c_{I2}-c_{I3}+c_Y+2c_{Y2}+2c_{YI1},
\]
and it vanishes, $K_{\rm pf}=0$, on the Solar-compatible coefficient family.
This is the local coefficient condition used here for
$\alpha_1=\alpha_2=\alpha_3=0$. A full vector-sector PPN analysis against the
bounds $|\alpha_1|\lesssim10^{-4}$, $|\alpha_2|\lesssim10^{-7}$
\cite{Will2014,WillNordtvedt1972,Kostelecky2011} belongs to the complete PPN
extension and is the standard follow-up for any theory with a solid sector.

The static light index $n_{\rm light}=\sqrt{A/B}$ expands as
\[
n_{\rm light}=1+2U+\frac74U^2+O(U^3)
\qquad(\gamma=1,\ \beta=1,\ b_2=\tfrac32),
\]
so the second-order optical coefficient is $q_{2{\rm PN}}=7/4$. These values
$\gamma=\beta=1$, $b_2=3/2$, and
$q_{2{\rm PN}}=7/4$ coincide with general relativity through second
post-Newtonian order. The weak-field sector of RefG is therefore a
consistency check --- measurements from Cassini,
$|\gamma-1|=(2.1\pm2.3)\times10^{-5}$ \cite{Bertotti2003}, and lunar laser
ranging constrain it exactly as they do GR --- rather than a discriminating
test; the fully coupled self-consistent exterior ODE and a raw Solar-System
likelihood remain separate steps. The Solar branch imposes the second-order
relation $c_{YI1}=2c_{Y2}$, whose kinetic consequence is handled in
Section~\ref{sec:stability}. The radial medium deformation that absorbs the
second-order angular residual is derived in Appendix~\ref{app:solar}.

Two scaling statements close the weak-field analysis. First, the technical
background
completion contributes $\Delta_{\rm Solar}\sim\Lambda R_\odot^2\sim
5.3\times10^{-35}$ to the Solar metric, which leaves the post-Newtonian
matching unchanged. Second, the projected deficit operator of the compact
branch is suppressed in this regime: relative to the Newtonian curvature
scale, its local
source ratio is $Ce^{-C}/4\approx1.06\times10^{-6}$ at the Solar surface
($C=r_s/R_\odot$). Consequently, the diffuse medium stress dominates the
weak-Solar matching, and the two branches do not contaminate each other.

Finally, the total localized energy fixes both normalizations used above:
the same $E_0$ serves as both the Noether inertial mass and the far-zone
source charge,
$M_{\rm inert}=M_{\rm source}=E_0/c^2$, in the standard minimally coupled
rest-energy normalization.

\section{Local stability of the scalar sector}
\label{sec:stability}

The Solar second-order condition $c_{YI1}=2c_{Y2}$ sets the bare longitudinal
solid kinetic prefactor of the minimal polynomial to zero,
$K_{\pi\pi}=2c_{Y2}-c_{YI1}=0$, while the phase kinetic term remains positive,
$K_{\Phi\Phi}=4c_{Y2}>0$. The physical Solar point therefore sits on a
kinetic degeneracy of $F_{\rm min}$ that must be lifted by the medium's
dynamic response. Earlier single-constraint completions overconstrain the
phase lag and the longitudinal response and pin the scalar speeds to a
defective luminal double root. That diagnostic, together with the bare
no-ghost window of the minimal polynomial, is collected in
Appendix~\ref{app:stability}.

The mechanism used to lift this degeneracy is a static-silent dynamic
operator. With
$U^A=u^\mu\partial_\mu\phi^A$, $Q^A=-g^{\mu\nu}\partial_\mu\phi^A
\partial_\nu\Phi$, and $W^A=U^A+Q^A$ (locally $W_L=\dot\pi_L$,
$Q_L=\nabla_L\chi$), the operator
\[
\Delta\mathcal L_{\rm dyn}
=\epsilon_B W_A W^A+2\epsilon_M W_A Q^A
\]
vanishes for any static source, so the Solar exterior of
Section~\ref{sec:weak} is untouched. In the quadratic action, however, it
shifts only the longitudinal kinetic entry and the mixing entry of the
scalar-longitudinal principal matrix,
\[
B_{\rm new}=c_Z+\epsilon_B,\qquad
M_{\rm new}=8c_{Y2}-4\lambda_6-c_Z+\epsilon_M .
\]
On the representative Solar slice $c_{Y2}=1$, $\lambda_6=1/4$, $c_Z=1$, the
completed determinant is
\[
\det M_{\rm rep}(s)
=5(1+\epsilon_B)s^2+\left[26+\epsilon_B-(6+\epsilon_M)^2\right]s+5,
\qquad s=\omega^2/k^2,
\]
and real, positive, strictly subluminal roots exist in the open window
\[
26+\epsilon_B+10\sqrt{1+\epsilon_B}<(6+\epsilon_M)^2<36+6\epsilon_B,
\]
of width $5(\sqrt{1+\epsilon_B}-1)^2$, finite for every $\epsilon_B>0$. At
the representative point $\epsilon_B=1$, $\epsilon_M=\sqrt{83/2}-6$, the roots
are
\[
s_1=\frac{29-\sqrt{41}}{40}\simeq0.5649,\qquad
s_2=\frac{29+\sqrt{41}}{40}\simeq0.8851,
\]
both of which are subluminal (Figure~\ref{fig:stability-window}). The
falsification object is the window itself, not the representative point: a
ghost, a complex root, or a superluminal scalar-longitudinal root rejects the
parameter branch. The statements of this section are local flat-background
principal-symbol results; establishing curved-background hyperbolicity and
completing the constraint-closure analysis are open problems, listed in
Section~\ref{sec:discussion}.

\begin{figure}[t]
\centering
\includegraphics[width=0.78\textwidth]{figures/refg_stability_window.pdf}
\caption{Completed scalar-longitudinal principal-symbol window. The shaded
region is the ghost-free subluminal interval
$26+\epsilon_B+10\sqrt{1+\epsilon_B}<M_2<36+6\epsilon_B$, where
$M_2=(6+\epsilon_M)^2$. The point $(\epsilon_B,M_2)=(1,41.5)$ gives
$s_1\simeq0.565$ and $s_2\simeq0.885$.}
\label{fig:stability-window}
\end{figure}

\section{The compact phase branch}
\label{sec:compact}

\subsection{Exponential exterior}
\label{sec:exterior}

The compact branch uses the reduced exterior static phase equation
$(r^2\phi')'=0$ with a flat normalization at infinity, $\phi=-r_s/r$, and the
biconformal map $h=-\phi/2$. Together these give the exterior
\[
ds^2=e^{-r_s/r}c^2dt^2
-e^{r_s/r}\left(dr^2+r^2d\Omega^2\right),
\qquad r_s=\frac{2GM}{c^2}.
\]
This is the exponential (Papapetrou) metric
\cite{Papapetrou1954,Yilmaz1958,Boonserm2018}. Because $AB=1$, the curved
static harmonic current reduces exactly to $-r^2h'$, so the reduced phase
equation is the full curved equation on this branch. For comparison with the
literature, which uses $m=GM/c^2$ and the form $e^{-2m/r}$, note that
$r_s=2m$ throughout. The metric has no finite-radius Killing horizon. It
matches the Schwarzschild metric through the first post-Newtonian order and,
in the time component, through $O(m^2/r^2)$, departing at second order in the
spatial sector. Its Einstein tensor has the mixed diagonal profile
\[
G^t{}_t=-D,\qquad
G^r{}_r=+D,\qquad
G^\theta{}_\theta=G^\phi{}_\phi=-D,
\qquad
D=\frac{r_s^2\,e^{-r_s/r}}{4r^4}.
\]
Any source of this exterior must reproduce this four-component pattern:
opposite signs in the radial and tangential directions, with a single common
profile function $D(r)$. In GR, the canonical realization is a massless ghost
scalar, and the radial null energy condition is violated at every radius
\cite{Boonserm2018}. The remainder of this section shows how the RefG scalar
sector supplies this source without a propagating ghost, identifies the
finite-core window that realizes it, and shows how a regular core excises the
wormhole throat at $r=r_s/2$ \cite{Boonserm2018}.

\subsection{The raw-invariant obstruction}
\label{sec:obstruction}

We first record why the polynomial sector itself cannot source this exterior
--- the obstruction that motivates the phase-normalized frame of
Section~\ref{sec:branches}. Evaluating the raw invariants on the compact
identity medium configuration gives $Y=w$, $\lambda_r=\lambda_t=1/w$ with
$w=e^{r_s/r}$. Inserting these into $F_{\rm min}$ on the physical Solar
coefficient slice produces the stress
\[
\frac{\Theta_F^t{}_t}{M_*^4c_{Y2}}
=\frac{(w-1)(3w^4-5w^3+w^2-23w+16)}{w^3},
\qquad
\frac{\Theta_F^r{}_r}{M_*^4c_{Y2}}
=-\frac{(w-1)(w^4-7w^3-5w^2+3w+16)}{w^3},
\]
with the tangential components equal to the radial one. For example,
$\Theta_F^r{}_r/(M_*^4c_{Y2})=19/4$ at $w=2$. No constant retuning of the
source can absorb this tensor: the raw radial and tangential stresses are
equal, while the exponential geometry requires them to have opposite signs.
The insertion of raw invariants is therefore excluded as a compact source,
and the compact branch must instead be --- and is --- formulated on the
phase-normalized action, where the same polynomial is exactly quiet:
$\widehat Y=\widehat\lambda_r=\widehat\lambda_t=1$, so
$\widehat F_{\rm min}=0$, $\widehat\Theta_F^\mu{}_\nu=0$, and the $\phi^A$
Euler residual vanishes. The obstruction thus fixes the role of the phase
normalization: it is the unique frame of the single action of
Section~\ref{sec:action} in which this tensor is absent, and its weak-field
limit ($H=0$) reproduces every Solar result of Section~\ref{sec:weak}
unchanged.

\subsection{The projected deficit source: realizability}
\label{sec:realizability}

The active compact source is a single operator of the higher EFT sector,
written with an independent deficit scalar $H_\Delta$:
\[
u_\mu=\frac{\partial_\mu\Phi}{\sqrt Y},\qquad
\gamma^{\mu\nu}=u^\mu u^\nu-g^{\mu\nu},\qquad
Z_\Delta^\perp=
\gamma^{\mu\nu}\partial_\mu H_\Delta\,\partial_\nu H_\Delta,\qquad
\mathcal L_\Delta^\perp=
\omega_\Delta\,\frac{Z_\Delta^\perp}{8\pi G},
\]
where $\omega_\Delta$ is a dimensionless coefficient set to $1$ by the
branch normalization below. On the weak-Solar branch this channel is
unloaded, and its local source ratio at the Solar surface is $\sim10^{-6}$
(Section~\ref{sec:solar}), so the branches do not mix. On the compact
exterior, $H_\Delta=h=r_s/(2r)$ solves its own Euler equation,
$(r^2H_\Delta')'=0$, with zero residual. The operator's stress is
\[
\Theta^t{}_t=-\Delta_P,\qquad
\Theta^r{}_r=+\Delta_P,\qquad
\Theta^\theta{}_\theta=\Theta^\phi{}_\phi=-\Delta_P,
\qquad
\Delta_P=\frac{Z_\Delta^\perp}{8\pi G}
=\frac{r_s^2\,e^{-r_s/r}}{32\pi G\,r^4}.
\]
Comparing with the Einstein tensor above, we find $D=8\pi G\,\Delta_P$ and
\[
G^\mu{}_\nu-8\pi G\,\Theta^\mu{}_\nu=0
\]
in all four diagonal components. This is the central result of the section:
the match is a four-component functional identity with one common profile,
holding at every radius. While a single fitted constant could rescale an
overall amplitude, it could not simultaneously produce the correct relative
signs and the common radial profile in all four components. The only quantity
fixed by hand is the overall normalization
$\omega_\Delta=1$.

Translated into effective-fluid language, the active source has
$\rho_a=-\Delta_P$, $p_{r,a}=-\Delta_P$, $p_{t,a}=+\Delta_P$. As a result,
the radial null energy condition is violated at every exterior radius,
$\rho_a+p_{r,a}=-2\Delta_P<0$, while the transverse one is saturated. This
reproduces, as it must, the source pattern that
Ref.~\cite{Boonserm2018} proved unavoidable for this geometry; the physical
reading in RefG is a phase-pressure deficit of the medium, with the largest
exterior deficit at the matching surface.

The difference from the phantom realization is structural. A ghost scalar
carries a propagating degree of freedom with wrong-sign time kinetic energy.
Here, the projector $\gamma^{\mu\nu}$ removes the time-kinetic direction of
$H_\Delta$ on the static branch: the operator contributes no
$\dot H_\Delta^2$ term, and its second variation is a positive
static stiffness,
\[
\delta^2\mathcal L_\Delta^\perp
=\frac{e^{-r_s/r}}{8\pi G}
\left[
(\delta h_r)^2
+\frac{(\delta h_\theta)^2}{r^2}
+\frac{(\delta h_\phi)^2}{r^2\sin^2\theta}
\right]>0,
\]
and $H_\Delta$ enters the static problem as an elliptic constraint field
rather than as a phantom wave. The propagating scalar spectrum of the theory
is carried by the $F_{\rm min}$ sector analyzed in
Section~\ref{sec:stability}. The null-energy violation of the effective
source is thus achieved without a wrong-sign kinetic term in any propagating
sector --- which is what a ghost-free realization of this geometry requires.
Determining the full perturbation spectrum around the compact background,
including the coupled $\{g,\Phi,\phi^A,H_\Delta\}$ system, remains an open
dynamical problem, recorded in Section~\ref{sec:discussion}.

\subsection{Branch selector and the parameter-free window}
\label{sec:selector}

Which objects realize this branch is decided by a finite-core load selector,
and the selector closes into a parameter-free statement. Define the load
$Q=r_s/r_c$, where $r_c$ is the core radius, and the input compactness
$Q_0=Qe^{-Q/2}$. Surface-clock matching equates the volume-deficit clock to
the compact surface lapse,
\[
e^{-q_v/3}=e^{-Q/2}\quad\Longrightarrow\quad q_v=\frac{3Q}{2},
\]
Because the equation $Q_0=Qe^{-Q/2}$ alone has two Lambert branches,
$Q_\pm=-2W_{0,-1}(-Q_0/2)$, the input $Q_0$ by itself does not select the
branch. Instead, the selection is made by the source amplitude through the
deficit feedback
\[
q_v=S\,e^{-\chi q_v},\qquad
q_\star=\frac{W(\chi S)}{\chi},\qquad
\left.\frac{d}{dq}\left(Se^{-\chi q}-q\right)\right|_{q_\star}
=-\chi q_\star-1<0,
\]
This yields a stable fixed point for any $S,\chi>0$. Since
$S(Q)=\tfrac{3Q}{2}e^{3\chi Q/2}$ is monotone, the amplitude selects a unique
load.

The exponent $\chi$ is not an adjustable parameter, for two reasons that we
state as separate facts. First, it cannot depend on the polynomial
coefficients: the phase-normalized $F_{\rm min}$ sector carries zero stress
on this branch (Section~\ref{sec:obstruction}), so no $c_i$ or
$\omega_\Delta$ dependence can enter the feedback; $\chi$ is the exponent of
the metric deficit-readout filter expressed in the volume variable
$q_v=3h$. Second, the readout channel is fixed by what gravitates: the
external source amplitude is the ADM/Komar mass, a volume-extensive charge,
read through the cubic filter
\[
e^{-3h}=e^{-q_v}\quad\Longrightarrow\quad\chi=1,
\]
whereas the intensive clock/rod filter $e^{-h}=e^{-q_v/3}$ (exponent $1/3$)
measures a single rate, not an integrated mass, and is not the gravitating
channel. We classify the two facts: the coefficient-independence of
$\chi$ is an action-level statement, whereas the volume-readout
identification is a physical normalization argument --- the unique exponent
consistent with the extensive (Komar) character of the gravitating charge ---
and not an action-level derivation. With $\chi=1$ the selector becomes
\[
Q_\star=\frac{2W(S)}{3},
\]
Enforcing the requirement that the throat radius $r_s/2$ lie inside the
core while the photon sphere lies outside, $1<Q_\star<2$, yields the
parameter-free window
\[
\frac32e^{3/2}<S<3e^{3},
\qquad
\frac{r_s}{2}<r_c<r_s,
\qquad
e^{-1/2}<C_0<\frac{2}{e},
\]
where $C_0=Q e^{-Q/2}$ on the principal branch,
$Q(C_0)=-2W_0(-C_0/2)$. As an explicit realizing class, the Newtonian $n=1$
polytrope $p=K\rho_c^2$ has $C_0=4K\rho_c/c^2$ and enters the window for
$e^{-1/2}<4K\rho_c/c^2<2/e$; a relativistic equation-of-state map is the
natural refinement. The object-specific value of $S$ belongs to the core
equation of state or collapse history. The selector says which loads
realize the branch, not which astrophysical objects reach them; that
dynamical question is taken up in Section~\ref{sec:discussion}.

\subsection{The regular core: completeness of the geometry}
\label{sec:core}

Considered in isolation, the exterior is geodesically incomplete and
contains the wormhole throat at $r=r_s/2$ \cite{Boonserm2018}. Both features
are removed by a
$C^2$-matched core. With $x=r/r_c$ and $q=r_s/r_c\in(1,2)$, the logarithmic
core profiles
\begin{align*}
\log A_-&=q\left(\tfrac{35}{8}x^2-\tfrac{21}{4}x^4+\tfrac{15}{8}x^6\right),\\
\log B_-&=-q+q\left(-\tfrac{11}{8}x^2+\tfrac{9}{4}x^4-\tfrac{7}{8}x^6\right)
\end{align*}
match the exterior $\log A_+=q/x$ and $\log B_+=-q/x$ at $x=1$ in value,
slope, and curvature, so the Israel junction stress vanishes identically.
The combination obeys the closed identity
\[
\log(AB)_{\rm core}=q\,(x^2-1)^3\in[-q,0],
\]
so $(AB)_{\rm core}$ is bounded in $[e^{-q},1]$. The resulting geometry has
the following properties, each an explicit computation:
\begin{enumerate}
\item \emph{No horizon:} $B=e^{\log B}>0$ everywhere, including $B(0)=e^{-q}$.
\item \emph{Regular center:} $A(0)=1$, $A'(0)=B'(0)=0$, areal radius
$R(0)=0$, with finite curvature
$R_0=\dfrac{177\,q}{4r_c^2}$,
$K_0=\dfrac{15063\,q^2}{16r_c^4}$ (Ricci and Kretschmann scalars).
\item \emph{Finite depth:} the proper radial distance from the matching
surface to the center is finite (e.g., $1.74\,r_c$ at $q=3/2$), whereas the
bare exterior continued to $r=0$ has infinite proper depth --- the infinite
throat is an artifact removed by the core.
\item \emph{Geodesic completeness:} a radial null geodesic reaches the
regular center at finite affine parameter, bounded by the $(AB)$ identity
above, and extends through it; outgoing geodesics reach infinity at infinite
affine parameter.
\item \emph{Bounded continuous source:} the effective source
$G^\mu{}_\nu/8\pi G$ of the cored metric is bounded on the whole core,
continuous at the junction --- where its radial null load equals the
exterior value $-2\Delta_P$ --- and isotropic at the center, with
$8\pi G\,r_c^2\,\rho_0=-\tfrac{105q}{4}$ and
$8\pi G\,r_c^2\,p_0=6q$.
\end{enumerate}
The geodesic markers below are therefore markers of a complete, regular,
horizonless geometry. The limit of this construction is equally
explicit: the core profile is a $C^2$ matching ansatz whose source is read
from the geometry, in the same effective-source sense as the exterior; a
core derived from a matter equation of state or from the medium action is
the open interior problem.

\subsection{Geodesic markers and the light-ring pair}
\label{sec:geodesics}

The test-particle structure of the exponential exterior is known
\cite{Boonserm2018,MakukovMychelkin2018}, and we quote it in the present
normalization ($r_s=2m$). The photon sphere is located at $r_{\rm ph}=r_s$,
and the critical impact parameter is
\[
b_c=e\,r_s
\qquad\Longrightarrow\qquad
\frac{b_c}{b_c^{\rm GR}}=\frac{2e}{3\sqrt3}=1.0463,
\]
a $+4.63\%$ shadow-scale excess over a Schwarzschild black hole of the same
mass. For massive particles, circular orbits give the ISCO from the condition
$r^2-3r_sr+r_s^2=0$,
\[
r_{\rm ISCO}=\frac{3+\sqrt5}{2}\,r_s,
\qquad
\frac{f_{\rm ISCO}}{f_{\rm ISCO}^{\rm GR}}
=\left[\frac{54\,e^{-2/\varphi^2}}{\varphi^4(2\varphi^2-1)}\right]^{1/2}
=0.931,
\]
a $6.9\%$ reduction of the ISCO frequency, in agreement with
\cite{MakukovMychelkin2018}. (The exterior root can also be written
$r_{\rm ISCO}=\varphi^{2}r_s$, with $\varphi$ the golden ratio.) What the present
construction adds to these known numbers is their setting: they are now
markers of a geodesically complete, regular geometry --- its exterior realized
by the projected EFT source, its core carrying a bounded, junction-continuous
effective source not yet derived from an equation of state --- in a
parameter-free compactness window.

The light-ring structure of the cored object follows from the optical
potential $V=B/(Ar^2)$ computed through the matched geometry. The exterior
maximum at $r=r_s$ is the unstable light ring. The potential $V$ is
continuous and $C^1$ at the junction, with $V'(r_c)>0$ for $q>1$, and rises
as the centrifugal wall $e^{-q}/r^2$ at the regular center. Consequently,
$V$ has exactly one interior minimum --- a stable light ring inside the core
--- realizing the light-ring pairing required for horizonless ultracompact
objects \cite{CunhaBertiHerdeiro2017}. Its location and the depth of the trapping
well are explicit: at $q=3/2$ the stable ring sits at $r\simeq0.87\,r_c$
($0.58\,r_s$), and the well-depth ratio between the unstable maximum and the
stable minimum is
\[
\frac{V_{\rm max}}{V_{\rm min}}
=1.06,\ 1.33,\ 2.05
\qquad\text{at}\qquad
q=1.2,\ 1.5,\ 1.9,
\]
i.e.\ a shallow well across the whole realizing window, shallowest at low
load. The dynamical implications are discussed in
Section~\ref{sec:discussion}.

\begin{figure}[t]
\centering
\includegraphics[width=\textwidth]{figures/refg_geodesic_structure.pdf}
\caption{Strong-field geodesic diagnostics in units $x=r/r_s$. Panel (a):
null optical potential $B_{\rm ph}=e^{-2r_s/r}/r^2$ for the exponential
exterior and $B_{\rm ph}^{\rm GR}=(1-r_s/r)/r^2$ for Schwarzschild, showing
$r_{\rm ph}=r_s$, $b_c=e\,r_s$, and $b_c/b_c^{\rm GR}=1.0463$. Panel (b):
massive effective potential with the circular-orbit and marginal-stability
conditions giving $r_{\rm ISCO}=(3+\sqrt5)r_s/2$ for the exponential
exterior and $3r_s$ for Schwarzschild, with
$f_{\rm ISCO}/f_{\rm ISCO}^{\rm GR}=0.931$.}
\label{fig:geodesic-structure}
\end{figure}

\section{Discussion: observations, stability, and open problems}
\label{sec:discussion}

\subsection{Confrontation with observations}
\label{sec:observations}

Table~\ref{tab:observations} summarizes the confrontation of both branches
with the data. The weak-field rows are GR-degenerate by construction; the
compact-branch rows are the discriminating content of the theory.

\begin{table}[t]
\centering
\small
\renewcommand{\arraystretch}{1.25}
\begin{tabular}{L{0.27\textwidth}L{0.14\textwidth}L{0.2\textwidth}L{0.3\textwidth}}
\toprule
Observable & GR & RefG & Bound / instrument \\
\midrule
PPN $\gamma-1$ & $0$ & $0$ &
$(2.1\pm2.3)\times10^{-5}$, Cassini \cite{Bertotti2003} \\
PPN $\beta-1$ & $0$ & $0$ & LLR/perihelion \cite{Will2014} \\
Preferred frame $\alpha_{1,2}$ & $0$ & $0$ (displayed order) &
$10^{-4}$, $10^{-7}$ \cite{Will2014} \\
2PN optical $q_{2\rm PN}$ & $7/4$ & $7/4$ & below current sensitivity \\
Tensor speed $c_g/c-1$ & $0$ & $0$ identically &
${\rm few}\times10^{-15}$ \cite{Abbott2017GRB} \\
\midrule
Shadow radius (same mass) & $\tfrac{3\sqrt3}{2}r_s$ & $e\,r_s$ ($+4.63\%$) &
Sgr~A*: allowed at $2\sigma$, $\sim1$--$1.5\sigma$ tension
\cite{EHT2022SgrA,Vagnozzi2023}; decidable by ngEHT \\
ISCO frequency (same mass) & $f^{\rm GR}_{\rm ISCO}$ &
$0.931\,f^{\rm GR}_{\rm ISCO}$ & X-ray QPO / continuum fitting \\
\bottomrule
\end{tabular}
\caption{Confrontation of RefG with observations. Weak-field entries
coincide with GR (consistency checks); compact-branch entries are
predictions of the static branch for objects in the realizing window.}
\label{tab:observations}
\end{table}

The shadow prediction deserves a precise statement. The Event Horizon
Telescope measures the fractional deviation of the Sgr~A* shadow diameter
from Schwarzschild as $\delta=-0.04^{+0.09}_{-0.10}$ (Keck mass-to-distance
prior) and $\delta=-0.08^{+0.09}_{-0.09}$ (VLTI prior) \cite{EHT2022SgrA}.
The static branch predicts $\delta=+0.046$. This value lies within the
$2\sigma$ bands but about $1$--$1.5\sigma$ above the measured central values
\cite{Vagnozzi2023}. Two caveats accompany any such comparison: the
ring-to-shadow calibration depends on the emission model, and the
prediction is for the static branch while Sgr~A* may spin. With those
caveats, this is the most directly testable consequence of the model: because
the window is parameter-free, the prediction cannot be shifted toward the
Schwarzschild value, and next-generation horizon-scale measurements can
falsify the branch.

\subsection{Light-ring stability and the gravitational-wave channel}
\label{sec:lightring}

Section~\ref{sec:geodesics} established that the cored object carries the
generic light-ring pair of horizonless ultracompact objects: the unstable
exterior ring at $r_s$ and one stable ring inside the core
\cite{CunhaBertiHerdeiro2017}. A stable light ring traps null modes, and
nonlinear evolutions of ultracompact boson and Proca stars show that the
associated pile-up can drive migration or collapse on timescales of order
$10^3$ light-crossing times \cite{CunhaEtAl2023}. Whether this fate applies
here is an open dynamical question, and we state the one quantitative fact
this paper contributes to it: the trapping well of the cored geometry is
shallow --- the potential ratio between the unstable maximum and the stable
minimum ranges from $1.06$ to $2.05$ across the realizing window, compared
with the deep cavities of the configurations evolved in
\cite{CunhaEtAl2023}. A shallow well stores little energy per trapped mode
and leaks it across a low barrier; whether that suffices for long-term
viability requires a dedicated nonlinear evolution and is left open here. In the gravitational-wave channel, the prompt ringdown of the cored
object is controlled by the same exterior light ring as in GR, while the
regular core replaces the horizon as the inner boundary, which is the
classic setting for late-time echoes \cite{Cardoso2016,AbediEtAl2017,
WesterweckEtAl2018}; the quasinormal spectrum of the cored geometry remains
open.

\subsection{Falsifiability and open problems}
\label{sec:open}

The falsification conditions for the branch and for the framework are
distinct, and we state both. The compact branch is excluded if
horizon-scale shadow measurements of an object independently shown to lie in
the realizing window exclude a $+4.6\%$ excess. It is similarly excluded if
the trapped-mode instability is shown to destroy the cored configuration on
astrophysical timescales, or if the rotating extension (once constructed)
fails the existing spin-sensitive constraints. The framework as a whole loses its
discriminating content with the branch: its weak-field sector is
GR-degenerate, so RefG without the compact branch is observationally
indistinguishable from GR at currently accessible orders. The
scalar-sector stability window of
Section~\ref{sec:stability} is an independent internal filter: a
demonstrated ghost or superluminal root on the physical coefficient family
rejects the parameter branch.

The open problems, each with its scope boundary, are: (i) the rotating
exterior --- all statements here are static; rotation is required before
spin-sensitive data can be confronted; (ii) the quasinormal-mode and echo
spectrum of the cored geometry, including the nonlinear fate of the
trapped modes; (iii) an interior derived from a matter equation of state or
from the medium action, replacing the $C^2$ matching ansatz whose effective
source we exhibited; (iv) the collapse dynamics that decides whether a
given astrophysical object reaches the realizing window --- the selector of
Section~\ref{sec:selector} is a static statement about loads, not a
formation theory; (v) curved-background hyperbolicity and the constraint
closure of the scalar sector, beyond the flat principal-symbol analysis;
(vi) the cosmological perturbation theory and the FLRW tensor mass term of
the solid sector; and (vii) the explicit vector-sector PPN parameters
$\alpha_1,\alpha_2$ beyond the displayed coefficient condition. None of
these affects the static results derived above; each is the standard next
layer of its respective literature, and (i)--(iii) are shared by every
horizonless-object program, including gravastars and boson stars
\cite{CardosoPani2019}.

\section{Conclusion}
\label{sec:conclusion}

RefG is a supersolid-class gravitational EFT with three established
properties and one parameter-free prediction. The weak static sector reproduces
Newton's law through a refractive source--index chain and coincides with GR
through the second post-Newtonian order; the tensor sector is luminal
identically; and the scalar-longitudinal sector has an open, subluminal stability
window at the local level. On the compact branch, the projected
phase-pressure deficit of the same scalar sector sources the exponential
(Papapetrou) exterior exactly --- a geometry whose GR realizations require a
phantom scalar --- through a four-component functional match with no
propagating ghost. The feedback exponent is fixed by the volume readout of
the asymptotic mass, making the realizing window $r_s/2<r_c<r_s$
parameter-free. A $C^2$-matched core renders the geometry horizonless,
regular at the center, and geodesically complete, while excising the wormhole
throat of the bare exterior. The branch inherits the known geodesic numbers
of the exponential metric and turns them into falsifiable predictions of a
complete geometry: a $+4.6\%$ shadow excess, already under mild tension from
Sagittarius~A* and decidable by next-generation instruments, and an ISCO
frequency $0.931$ of the GR value at the same mass. The rotating exterior, the
quasinormal spectrum with the trapped-mode question, the equation-of-state
interior, and the collapse dynamics of branch selection are the next layer
of the program.

\section*{Data availability statement}

No observational or experimental datasets were generated or analyzed in
this study. The results are analytic; the algebraic identities and branch
gates underlying every displayed equation are implemented in an
accompanying Python/SymPy source tree (sector gates p01--p16), publicly
archived at \texttt{doi:10.5281/zenodo.20645891}. A README in the archive
maps each displayed equation to its gate; in particular, the source identity
of Section~\ref{sec:realizability} is verified in
\texttt{p05g}/\texttt{p05i}, the raw-invariant obstruction in
\texttt{p05j}/\texttt{p05k}, the selector and $\chi=1$ in
\texttt{p05y}/\texttt{p16g}/\texttt{p16h}, the cored geometry and light
rings in \texttt{p16j}--\texttt{p16l}, the stability window in
\texttt{p01}, and the Solar sector in \texttt{p03}--\texttt{p03d}.

\appendix

\section{Stability diagnostics of the minimal sector}
\label{app:stability}

This appendix records the two diagnostic stages that precede the completed
window of Section~\ref{sec:stability}. Around the homogeneous background,
with $\Sigma=c_Y+3c_{YI1}$, the bare phase and longitudinal kinetic entries
of the minimal polynomial are
\[
K_{\Phi\Phi}=\Sigma+6c_{Y2},\qquad
K_{\pi\pi}=-\Sigma-2c_{Y2},
\]
so the nondegenerate bare no-ghost window is $c_{Y2}>0$,
$-6c_{Y2}<\Sigma<-2c_{Y2}$. The Solar-compatible coefficient family is
\begin{align*}
c_{I1}&=4c_{Y2}+2c_{YI1}, &
c_{I1sq}&=c_{Y2}, &
c_{I2}&=-10c_{Y2}-3c_{YI1},\\
c_{I3}&=8c_{Y2}+4c_{YI1}, &
c_Y&=-4c_{Y2}-2c_{YI1}, & &
\end{align*}
and the physical Solar $2$PN condition $c_{YI1}=2c_{Y2}$ puts the
longitudinal entry on the boundary, $K_{\pi\pi}=0$, with
$K_{\Phi\Phi}=4c_{Y2}>0$ --- the degeneracy lifted in
Section~\ref{sec:stability}.

An earlier background completion used the combination $C_6=Y+2I_1-I_3$ and
the invariant $Z=\delta_{AB}(u^\mu\partial_\mu\phi^A)
(u^\nu\partial_\nu\phi^B)$ through the quadratic block
$\mathcal L_2=\lambda_6(\delta C_6)^2+c_ZZ$ with
$\delta C_6=2(\dot\chi+\partial_i\pi^i)$ and
$Z=(\dot\pi_i-\partial_i\chi)^2$. Its scalar-longitudinal determinant on the
Solar slice has Fourier coefficients
\[
A=D=4(c_{Y2}+\lambda_6),\qquad B=C=c_Z,\qquad
M=8c_{Y2}-4\lambda_6-c_Z,
\]
hence $p_0=p_2$ in $\det M_\odot(s)=p_2s^2+p_1s+p_0$: the product of the two
roots is one, so two distinct real roots cannot both be subluminal. At the
representative point $c_{Y2}=1$, $\lambda_6=1/4$, $c_Z=1$ one finds
$\det M_\odot(s)=5(s-1)^2$ --- a defective luminal double root. This block
therefore overconstrains the phase lag and the longitudinal response in a
single combination; the static-silent operator $\Delta\mathcal L_{\rm dyn}$
of Section~\ref{sec:stability} separates the two channels and opens the
subluminal window.

\section{Solar second-order derivations}
\label{app:solar}

\subsection*{Medium strain and $a_2=4$}

At second order the angular field equation of the frozen-solid ansatz leaves
a residual proportional to $(2c_{Y2}+c_{YI1})$. Allowing the radial medium
deformation $\phi^r=r(1+s(r)\,u^2+\cdots)$ converts this residual into the
balance equation
\[
2rs'(r)+2s(r)+1=0
\qquad\Longrightarrow\qquad
s(r)=\frac{C}{r}-\frac12 ,
\]
whose constant part fixes the local second-order coefficient; the $C/r$ term
is a boundary tail. With this deformation the metric part of the solution is
the GR-compatible isotropic $2$PN exterior, corresponding to $a_2=4$ in the
areal chart of Section~\ref{sec:solar}, and hence $\beta_A=1$. The bare
longitudinal degeneracy introduced by the same condition,
$K_{\pi\pi}=0$, is lifted by the static-silent effective-solid operator
$\delta_{AB}(u^\mu\partial_\mu\phi^A)(u^\nu\partial_\nu\phi^B)$, which
vanishes on the background and in the first static variation, leaving the
equation for $s(r)$, the value $a_2=4$, and $q_{2{\rm PN}}=7/4$ unchanged,
while contributing positive solid kinetic energy at perturbation level.

\subsection*{The $H$ constraint on the Solar branch}

Writing the weak-Solar strain as $Y=1+y_1U+y_2U^2$,
$\lambda_r=1+r_1U+r_2U^2$, $\lambda_t=1+t_1U+t_2U^2$, the auxiliary
constraint of Section~\ref{sec:branches} expands as
\[
\frac{S_H}{c_{Y2}}
=8(y_1+r_1+2t_1)\,U
+4\left(r_1^2+20r_1t_1+2r_2+12t_1^2+4t_2-y_1^2+2y_2\right)U^2
+O(U^3).
\]
The Solar strain has $y_1=1$, $r_1=-1$, $t_1=0$, which annihilates the first
order; the exact-GR $2$PN branch has $y_2=1$, $r_2=1$, $t_2=-1$, which
annihilates the second. Thus $H=0$ is an $H$-equation solution through the
displayed Solar order, not a formal reduction; the first-order cancellation
follows from the weak variation of $Y+I_1$ and the second-order one from the
same exact-GR vacuum condition that removes the $F_{\rm min}$ gradient. On
the compact pure-phase branch the hatted invariants equal one and $S_H=0$ is
exact.

\subsection*{Rejection of the determinant lock}

The trial identification $H=-\ln(\lambda_r\lambda_t^2)/6$, which would tie
the normalization to the solid determinant, gives at first Solar order
\[
\Theta_F^t{}_t=M_*^4c_{Y2}\,\frac{8}{3}\,U,\qquad
\Theta_F^r{}_r=\Theta_F^\theta{}_\theta=M_*^4c_{Y2}\,\frac{8}{9}\,U,
\]
a nonzero weak-field stress, and admits no exact-GR constant-strain $2$PN
solution. The determinant lock is therefore rejected, and $H$ is the
independent pressure/energy-deficit normalization used in the text.

% bibliography

\begin{thebibliography}{99}

\bibitem{Will2014}
C. M. Will, The confrontation between general relativity and experiment,
\textit{Living Reviews in Relativity} \textbf{17}, 4 (2014).

\bibitem{Bertotti2003}
B. Bertotti, L. Iess, and P. Tortora, A test of general relativity using
radio links with the Cassini spacecraft, \textit{Nature} \textbf{425},
374--376 (2003).

\bibitem{Abbott2017GW}
B. P. Abbott et al. (LIGO Scientific and Virgo Collaborations), GW170817:
Observation of gravitational waves from a binary neutron star inspiral,
\textit{Physical Review Letters} \textbf{119}, 161101 (2017).

\bibitem{Abbott2017GRB}
B. P. Abbott et al., Gravitational waves and gamma-rays from a binary
neutron star merger: GW170817 and GRB 170817A, \textit{Astrophysical
Journal Letters} \textbf{848}, L13 (2017).

\bibitem{EHT2019M87}
Event Horizon Telescope Collaboration, First M87 Event Horizon Telescope
results. I. The shadow of the supermassive black hole,
\textit{Astrophysical Journal Letters} \textbf{875}, L1 (2019).

\bibitem{EHT2022SgrA}
Event Horizon Telescope Collaboration, First Sagittarius A* Event Horizon
Telescope results. VI. Testing the black hole metric,
\textit{Astrophysical Journal Letters} \textbf{930}, L17 (2022).

\bibitem{CardosoPani2019}
V. Cardoso and P. Pani, Testing the nature of dark compact objects: a
status report, \textit{Living Reviews in Relativity} \textbf{22}, 4 (2019).

\bibitem{CliftonEtAl2012}
T. Clifton, P. G. Ferreira, A. Padilla, and C. Skordis, Modified gravity
and cosmology, \textit{Physics Reports} \textbf{513}, 1--189 (2012).

\bibitem{NicolisEtAl2015}
A. Nicolis, R. Penco, F. Piazza, and R. Rattazzi, Zoology of condensed
matter: framids, ordinary stuff, extra-ordinary stuff, \textit{Journal of
High Energy Physics} \textbf{06} (2015) 155.

\bibitem{CeloriaComelliPilo2017}
M. Celoria, D. Comelli, and L. Pilo, Fluids, superfluids and supersolids:
dynamics and cosmology of self-gravitating media, \textit{Journal of
Cosmology and Astroparticle Physics} \textbf{09} (2017) 036.

\bibitem{Endlich2013}
S. Endlich, A. Nicolis, and J. Wang, Solid inflation, \textit{Journal of
Cosmology and Astroparticle Physics} \textbf{10} (2013) 011.

\bibitem{Cheung2008}
C. Cheung, P. Creminelli, A. L. Fitzpatrick, J. Kaplan, and L. Senatore,
The effective field theory of inflation, \textit{Journal of High Energy
Physics} \textbf{03} (2008) 014.

\bibitem{Gubitosi2013}
G. Gubitosi, F. Piazza, and F. Vernizzi, The effective field theory of dark
energy, \textit{Journal of Cosmology and Astroparticle Physics}
\textbf{02} (2013) 032.

\bibitem{BelliniSawicki2014}
E. Bellini and I. Sawicki, Maximal freedom at minimum cost: linear
large-scale structure in general modifications of gravity, \textit{Journal
of Cosmology and Astroparticle Physics} \textbf{07} (2014) 050.

\bibitem{Horndeski1974}
G. W. Horndeski, Second-order scalar-tensor field equations in a
four-dimensional space, \textit{International Journal of Theoretical
Physics} \textbf{10}, 363--384 (1974).

\bibitem{Creminelli2017}
P. Creminelli and F. Vernizzi, Dark energy after GW170817 and GRB170817A,
\textit{Physical Review Letters} \textbf{119}, 251302 (2017).

\bibitem{Papapetrou1954}
A. Papapetrou, Eine Theorie des Gravitationsfeldes mit einer Feldfunktion,
\textit{Zeitschrift f\"ur Physik} \textbf{139}, 518--532 (1954).

\bibitem{Yilmaz1958}
H. Yilmaz, New approach to general relativity, \textit{Physical Review}
\textbf{111}, 1417--1426 (1958).

\bibitem{Misner1999}
C. W. Misner, Yilmaz cancels Newton, \textit{Il Nuovo Cimento B}
\textbf{114}, 1079 (1999), arXiv:gr-qc/9504050.

\bibitem{AlleyYilmaz1995}
C. O. Alley and H. Yilmaz, Refutation of C. W. Misner's claims,
\textit{Il Nuovo Cimento B} \textbf{114}, 1087 (1999),
arXiv:gr-qc/9506082.

\bibitem{Boonserm2018}
P. Boonserm, T. Ngampitipan, A. Simpson, and M. Visser, Exponential metric
represents a traversable wormhole, \textit{Physical Review D} \textbf{98},
084048 (2018).

\bibitem{MakukovMychelkin2018}
M. A. Makukov and E. G. Mychelkin, \textit{Physical Review D} \textbf{98},
064050 (2018).

\bibitem{Vagnozzi2023}
S. Vagnozzi et al., Horizon-scale tests of gravity theories and fundamental
physics from the Event Horizon Telescope image of Sagittarius A*,
\textit{Classical and Quantum Gravity} \textbf{40}, 165007 (2023).

\bibitem{WillNordtvedt1972}
C. M. Will and K. Nordtvedt, Conservation laws and preferred frames in
relativistic gravity. I. Preferred-frame theories and an extended PPN
formalism, \textit{Astrophysical Journal} \textbf{177}, 757--774 (1972).

\bibitem{Kostelecky2011}
V. A. Kosteleck\'y and N. Russell, Data tables for Lorentz and CPT
violation, \textit{Reviews of Modern Physics} \textbf{83}, 11--31 (2011),
arXiv:0801.0287.

\bibitem{CunhaBertiHerdeiro2017}
P. V. P. Cunha, E. Berti, and C. A. R. Herdeiro, Light-ring stability for
ultracompact objects, \textit{Physical Review Letters} \textbf{119},
251102 (2017).

\bibitem{CunhaEtAl2023}
P. V. P. Cunha, C. Herdeiro, E. Radu, and N. Sanchis-Gual, The fate of the
light-ring instability, \textit{Physical Review Letters} \textbf{130},
061401 (2023).

\bibitem{Cardoso2016}
V. Cardoso, S. Hopper, C. F. B. Macedo, C. Palenzuela, and P. Pani,
Gravitational-wave signatures of exotic compact objects and of quantum
corrections at the horizon scale, \textit{Physical Review D} \textbf{94},
084031 (2016).

\bibitem{AbediEtAl2017}
J. Abedi, H. Dykaar, and N. Afshordi, Echoes from the abyss: tentative
evidence for Planck-scale structure at black hole horizons,
\textit{Physical Review D} \textbf{96}, 082004 (2017).

\bibitem{WesterweckEtAl2018}
J. Westerweck et al., Low significance of evidence for black hole echoes
in gravitational wave data, \textit{Physical Review D} \textbf{97},
124037 (2018).

\end{thebibliography}

\end{document}
